\worksheettitle
  {Точки через равные промежутки}
  {\modulemark{M07} Учебный модуль \textbullet\ полный маршрут}

\studentnameblock

\begin{notebox}[Чему вы научитесь]
  Считать промежутки между позициями, находить расстояние, длину одного
  промежутка и номер конечной точки.
\end{notebox}

\begin{checkpointbox}[Вход: точки и переходы]
  Пять точек отмечены подряд через одинаковые расстояния.
  \[
    \text{точек } \answerblank[15mm],\qquad
    \text{промежутков } \answerblank[15mm].
  \]
  Почему числа различаются? \answerblank[70mm]
\end{checkpointbox}

\begin{notebox}[Как устроен маршрут]
  Сначала научитесь считать переходы по рисунку и номерам. Затем разберите
  прямую и обратную задачи, примените модель к ряду столбов, исправьте
  типичную ошибку и пройдите контроль.
\end{notebox}

\section{Точки и промежутки}

\pointsvsintervalsfigure

\begin{fillbox}[Разберите пример]
  Между первой и пятой точками $5-1=4$ промежутка. Поэтому пять подряд
  отмеченных точек образуют четыре перехода.
\end{fillbox}

\begin{notebox}[Правило]
  Число промежутков между точками равно разности их номеров:
  \[
    \text{промежутки}=\text{больший номер}-\text{меньший номер}.
  \]
\end{notebox}

\section{Практика с опорой}

\begin{practicebox}[1. Посчитайте промежутки]
  \begin{enumerate}[label=\alph*)]
    \item между 3-й и 9-й точками: $9-3=\answerblank[15mm]$;
    \item между 8-й и 27-й: $27-8=\answerblank[15mm]$;
    \item между 15-й и 16-й: \answerblank[15mm];
    \item между 41-й и 56-й: \answerblank[15mm].
  \end{enumerate}
\end{practicebox}

\studentcountstrip

\begin{practicebox}[2. Прочитайте рисунок]
  Между 2-й и 7-й точками \answerblank[15mm] промежутков.
  Если один промежуток равен $3$ см, расстояние равно
  \[
    \answerblank[15mm]\cdot3=\answerblank[18mm]\text{ см}.
  \]
\end{practicebox}

\section{Находим расстояние}

\distanceexamplefigure

\begin{notebox}[Алгоритм прямой задачи]
  \textbf{1.} Найдите число промежутков разностью номеров.
  \textbf{2.} Умножьте число промежутков на длину одного.
  \textbf{3.} Запишите единицу и проверьте схему.
\end{notebox}

\begin{practicebox}[3. Равноудалённые точки]
  Точки отмечены через $7$ мм. Найдите расстояние:
  \begin{enumerate}[label=\alph*)]
    \item между 5-й и 12-й: \answerblank[70mm];
    \item между 8-й и 27-й: \answerblank[70mm];
    \item между 31-й и 44-й: \answerblank[70mm].
  \end{enumerate}
\end{practicebox}

\section{Обратные задачи}

\reverseexamplefigure

\Needspace{7\baselineskip}
\begin{fillbox}[Разберите обратный ход]
  Между 3-й и 7-й точками $7-3=4$ промежутка. Общее расстояние $24$ см,
  поэтому один промежуток равен $24:4=\answerblank[15mm]$ см.
\end{fillbox}

\unknownpointfigure

\begin{practicebox}[4. Найдите номер точки]
  Точки расположены через $4$ см.
  \begin{enumerate}[label=\alph*)]
    \item От 12-й точки вправо прошли $28$ см. Номер:
      $12+28:4=\answerblank[18mm]$.
    \item От 30-й точки в сторону уменьшения номеров прошли $36$ см.
      Номер: $30-36:4=\answerblank[18mm]$.
  \end{enumerate}
\end{practicebox}

\section{Самостоятельная практика}

\postrowfigure

\begin{practicebox}[5. Столбы вдоль дороги]
  Семь столбов стоят через $4$ м, причём первый и последний входят в ряд.
  \begin{enumerate}[label=\alph*)]
    \item Число промежутков: \answerblank[18mm].
    \item Расстояние между крайними столбами: \answerblank[35mm].
  \end{enumerate}
\end{practicebox}

\begin{practicebox}[6. Обратная задача]
  От первого до последнего фонаря $54$ м. Фонари стоят через $6$ м и есть
  в обоих концах ряда. Сколько всего фонарей?
  \writinglines{3}
\end{practicebox}

\section{Разбираем ошибку}

\begin{warningbox}[«Десять точек --- десять промежутков»]
  Ученик решил, что между 1-й и 10-й точками $10$ промежутков.

  \textbf{Причина ошибки:} \answerblank[100mm]

  \textbf{Исправление:} \answerblank[100mm]
\end{warningbox}

\section{Контрольная точка}

\begin{checkpointbox}[7. Проверьте результат]
  \begin{enumerate}[label=\textbf{\arabic*.}]
    \item Точки отмечены через $5$ см. Найдите расстояние между 7-й и
      19-й точками.
    \item Расстояние между 4-й и 10-й точками равно $42$ мм. Найдите длину
      одного промежутка.
    \item На участке длиной $32$ м стоят столбы через $4$ м, включая оба
      конца. Сколько столбов?
    \item Объясните, почему объектов на один больше, чем промежутков.
  \end{enumerate}
  \writinglines{4}
\end{checkpointbox}

\begin{reflectionbox}[Проверяю себя]
  \begin{itemize}[label=\choicebox]
    \item различаю точки и промежутки;
    \item нахожу число переходов разностью номеров;
    \item решаю прямую и обратную задачу;
    \item учитываю, включены ли крайние объекты.
  \end{itemize}
\end{reflectionbox}

\begin{resultbox}[Дальнейший маршрут]
  Если считаете объекты вместо промежутков, выполните M07-R. Для тренировки
  используйте M07-H. После устойчивого результата продолжайте следующий
  раздел курса. M07\(\star\) выполняется дополнительно.
\end{resultbox}
